\section{From coordinates to distinct pairs}\label{sec:count}
For each \(d\) fix \(\rho_d\in R\) of norm \(d\), obtained by
multiplying one factor above each split prime of \(d\).
Embed the residual shell at \(m=n/d\) in the original shell
by \((z,y)\mapsto(\rho_d z,\rho_d y)\). This preserves the
fractional norm coordinate, the prescribed orientation at \(3\),
and inert divisibility above \(Z\).

Each original value pair recovers its triple \(d,u,v\).
At split primes up to \(Y\), the only possible positive valuation
is one, occurring in both norms at the primes of \(d\), only
in the first at \(u\), and only in the second at \(v\).
The images for distinct triples are consequently disjoint.
This makes it possible to bound the remaining losses on the
original shell without summing an error over the triples.

\subsection{Large inert primes}
If an inert prime \(p>Z\) divides the first component norm,
that component is \(p\beta\). The norm of \(\beta\) belongs to
\([(x-H/2)/p^2,(x+H/2)/p^2]\). For the other component use
\(n-x\) in place of \(x\). Each resulting annulus has
squared-radius width \(H/p^2\) and outer radius at most
\(\sqrt n/p\). Covering each lattice point by a translated
fundamental cell, of fixed area and bounded radius, gives
\(O(H/p^2+\sqrt n/p+1)\) possible \(\beta\). Indeed the cells
lie in the annulus obtained by moving the two boundaries by
a fixed distance, truncating the inner radius at zero.
The complementary positive norm has at most
\(6\max_{a\le n}\tau(a)=n^{o(1)}\) coordinate choices.

Since all weights are at most one, summing on the original
shell over the two components and \(Z<p\le\sqrt n\) bounds
the lost weight by
\[
 n^{o(1)}\left(\frac H Z+\sqrt n(1+\log n)\right)
             =o_\varepsilon(H/\log n).
\]
Here \(Z\) is a fixed power of \(n\), and the hypothesis
\(H^3\ge n^{1+\varepsilon}M\) with \(M\ge H\) gives
\(H\ge n^{(1+\varepsilon)/2}\). Thus \eqref{eq:weight}
persists after removing every inert prime. Oddness and the
condition at \(3\) now give the required norm values.

\subsection{The weight of one value pair}
On the residual shell write \(A_1=uk_1\) and \(A_2=vk_2\).
All split prime factors of \(k_1,k_2\) exceed \(Z=n^\eta\).
Each of these two numbers has at most \(1/\eta\) split factors
counted with multiplicity, and the possible single factor \(3\)
has no effect on coordinate multiplicity.
The split primes in \(uv\) occur to valuation one. By the
multiplicity formula in Section~\ref{sec:angular}, the weighted
coordinate multiplicity, after division by \(\tau(uv)\),
is at most
\[
 36f(k_1)f(k_2)\le36\cdot 2^{2/\eta}.
\]
Each original pair determines its triple and residual pair.
Dividing the surviving weight by this constant proves the lower
bound of Theorem~\ref{thm:main}.

\subsection{The upper bound on integers}
For the upper bound, write \(s=s_1\) in the interval of length
\(H\). At every inert \(p\ge5\) we must exclude the residues
\(0,n\bmod p\). Let \(\nu(p)\) be their number, namely one if
\(p\mid n\) and two otherwise, and extend \(\nu\)
multiplicatively to squarefree products of these primes.
The Chinese remainder theorem and the count in a residue
class give \(H\nu(\ell)/\ell+O(\nu(\ell))\) integers in
all the excluded sets for \(p\mid\ell\), uniformly in the
interval location and endpoint convention.

Use the upper polynomials of Section~\ref{sec:sieve} with
probabilities \(\nu(p)/p\), sieving to \(Z'=n^{c_1}\) for a
sufficiently small fixed \(c_1>0\).
The block probability bounds hold as before. Choose \(c_1\)
so that the corresponding divisor level is \(D'\le n^{1/4}\).
The total error is \(O(D'\log(2D'))=o(H/\log n)\), since
\(\nu(\ell)\le\tau(\ell)\)
and \(\sum_{\ell\le D'}\tau(\ell)\ll D'\log(2D')\).
The upper main term is at most a constant times \(H\) times
\[
 \prod_{\substack{5\le p\le Z'\\ p\ {\rm inert}}}(1-2/p)
 \prod_{\substack{5\le p\le Z',\ p\mid n\\ p\ {\rm inert}}}
                        \frac{p-1}{p-2}
 \ll \frac{J(n)}{\log n}.
\]
This proves the upper bound and completes Theorem~\ref{thm:main}.
The argument for this upper bound in fact works in any interval
whose length is at least a fixed positive multiple of \(\sqrt n\).

\begin{proof}[Proof of Corollary~\ref{cor:ranges}]
In a fixed interval of shares, take a smaller fixed interval
of length at most the distance of its center from the nearest
endpoint of \([0,1]\). This satisfies the hypotheses with
\(H\asymp n\) for large \(n\), giving the lower bound. Use the
upper sieve on the full interval and on \([0,n]\) for the
unrestricted count.
At \(x=\lambda n\) with \(0<\lambda<1\), a width \(H=n^\theta\)
satisfies the condition for any fixed
\(0<\varepsilon<3\theta-2\).
Finally, inside \([a n^\theta,b n^\theta]\) take a subinterval
of width comparable to \(n^\theta\) and at most its midpoint.
Its distance to the nearer endpoint of \([0,n]\) is comparable
to \(n^\theta\), so one can use \(0<\varepsilon<2\theta-1\).
The upper sieve on the full interval gives the matching bound,
also for \(a=0\).
\end{proof}
